% !TEX program = xelatex
% 編譯： xelatex 報告_VPP-RTS.tex  (跑兩次以產生目錄)
% 需 XeLaTeX + ctex；中文字型使用 PingFang TC（macOS）。
\documentclass[12pt,a4paper,fontset=none]{ctexart}

% ---- 中文字型（macOS 系統字型，繁體）----
\setCJKmainfont{PingFang TC}
\setCJKsansfont{PingFang TC}
\setCJKmonofont{PingFang TC}

% ---- 套件 ----
\usepackage{amsmath,amssymb}
\usepackage{geometry}
\geometry{margin=2.4cm}
\usepackage{booktabs}
\usepackage{array}
\usepackage{tabularx}
\usepackage{xcolor}
\usepackage{enumitem}
\usepackage{fancyvrb}
\usepackage{titlesec}
\usepackage[hidelinks,unicode]{hyperref}

\definecolor{vblue}{RGB}{31,73,125}
\definecolor{vgreen}{RGB}{46,125,50}
\definecolor{vgray}{RGB}{242,242,242}

\setlist[itemize]{itemsep=2pt,topsep=2pt,parsep=0pt}
\setlist[enumerate]{itemsep=2pt,topsep=2pt,parsep=0pt}

% 章節標題上色
\titleformat{\section}{\Large\bfseries\color{vblue}}{\thesection}{0.6em}{}
\titleformat{\subsection}{\large\bfseries\color{vblue}}{\thesubsection}{0.6em}{}
\titleformat{\subsubsection}{\normalsize\bfseries}{\thesubsubsection}{0.6em}{}

% 識別字（檔名 / 參數）：\fn{...}，內含底線請寫 \_
\newcommand{\fn}[1]{\texttt{\small #1}}

% 程式碼框
\DefineVerbatimEnvironment{code}{Verbatim}%
  {frame=single,framesep=4pt,fontsize=\small,rulecolor=\color{gray!60},samepage=false}

\newcommand{\R}{\mathbb{R}}
\renewcommand{\arraystretch}{1.25}
\renewcommand{\contentsname}{目錄}   % 繁體（ctex 預設為簡體「目录」）

\begin{document}

% =========================================================== 封面
\begin{titlepage}
\centering
\vspace*{2.5cm}
{\Huge\bfseries\color{vblue} 虛擬電廠即時排程系統\par}
\vspace{0.4cm}
{\LARGE\bfseries Virtual Power Plant Real-Time Scheduling (VPP-RTS)\par}
\vspace{1.2cm}
{\Large 即時系統 (Real-Time Systems) 作業\par}
\vspace{0.3cm}
{\large 完整實作報告　+　報告講稿\par}
\vspace{1.5cm}
{\large Level 1：日前排程 $\cdot$ 接收測試 $\cdot$ 效能評估\par}
{\large Level 2：放寬假設建模 $\cdot$ 進階動態排程\par}
\vspace{2.0cm}
{\large 國立成功大學　即時系統課程\par}
\vfill
{\small 本報告所有數據取自專案 \fn{output/} 與 \fn{input/} 之實際 JSON 輸出（固定隨機種子、可重現）。\par}
\end{titlepage}

% =========================================================== 文件用途
\section*{本文件用途}
\addcontentsline{toc}{section}{本文件用途}
本報告分為兩大部分：
\begin{itemize}
  \item \textbf{第一部分「完整實作說明」}：讓未參與開發的組員，讀完即可掌握整份作業的問題定義、數學模型、四個階段的演算法、Level 2 的放寬假設與動態排程，以及實際執行結果。
  \item \textbf{第二部分「報告講稿」}：依投影片順序撰寫，可直接照著對老師口頭報告（含開場、各段重點、預期問答）。
\end{itemize}

\tableofcontents
\newpage

% =========================================================== 第一部分
\part*{第一部分　完整實作說明}
\addcontentsline{toc}{section}{第一部分　完整實作說明}

% ---- 1 背景 ----
\section{作業背景與目標}
「虛擬電廠（Virtual Power Plant, VPP）」把分散的發電與儲能資產整合成一個可被統一調度的電力系統。本作業要為一個 VPP 設計一套即時排程系統：一方面安排各種「發電設備」供電，另一方面滿足一連串有時間限制的「用電需求（job）」。

我們把發電設備視為即時系統中的\textbf{處理器（processor）}、把用電需求視為\textbf{任務（task/job）}，於是電力調度問題就被轉化成一個即時排程問題。時間軸為 $H=72$ 小時，每格 $\Delta t = 1$ 小時。

\subsection{三類發電設備（處理器）}
\begin{itemize}
  \item \textbf{傳統機組 $I_g$}：可控輸出的火力機組，有開關機、最小開／關機時間、爬升率與成本。
  \item \textbf{再生能源 $I_r$}：太陽能，輸出受日前預測曲線上限約束。
  \item \textbf{儲能裝置 $I_b$}：電池，可充可放但不可同時，受 SOC 上下限約束。
\end{itemize}

\subsection{三類用電需求（任務）}
\begin{itemize}
  \item \textbf{週期任務 $J_p$}：固定週期釋放的硬截止任務（例：工廠定時製程）。
  \item \textbf{零星任務 $J_s$}：硬截止、臨時出現，須通過「接收測試」才接受（例：緊急冷卻）。
  \item \textbf{非週期任務 $J_a$}：軟截止、可等待，逾時僅記為 miss 與 tardiness（例：非緊急充電）。
\end{itemize}

\subsection{整體目標函數（三個彼此衝突的目標）}
系統以最小化形式同時追求「少漏非週期任務、低機組成本、高售電收益」：
\begin{equation}
\min\; F = \alpha f_1 + f_2 + f_3
\end{equation}
\begin{align}
f_1 &= \sum_{j\in J_a}\mathrm{Miss}_j, & &\alpha = 10000\ (\text{\$/miss}) \\
f_2 &= \sum_{i\in I_g}\sum_{t\in T}\Big(\mathrm{coston}_i\cdot\min(1,P_{i,t}) + \mathrm{costup}_i\cdot P_{i,t}\Big) \\
f_3 &= -\sum_{t\in T}\lambda_t\,\mathrm{Sell}_t
\end{align}
其中 $f_1$ 為非週期逾時數（以罰係數 $\alpha$ 轉成金額），$f_2$ 為機組成本，$f_3$ 為售電收益取負號（最小化 $f_3$ 即最大化收益）。

% ---- 2 輸入 ----
\section{系統輸入與資產設定}
輸入檔 \fn{input/processor\_settings.json} 與 \fn{input/price\_72hr.json} 定義固定的資產與價格。本次使用的資產如下。

\subsection{傳統機組（2 部）}
\begin{center}
\begin{tabularx}{\linewidth}{lccccc}
\toprule
\textbf{機組} & \textbf{出力下/上限(MW)} & \textbf{爬升 RU/RD} & \textbf{最小開/關(h)} & \textbf{固定(\$/h)} & \textbf{變動(\$/MWh)} \\
\midrule
thermal\_1 & 15 / 80 & 15 / 15 & 3 / 2 & 1200 & 42 \\
thermal\_2 & 10 / 45 & 20 / 20 & 2 / 2 & 600 & 70 \\
\bottomrule
\end{tabularx}
\end{center}
\noindent\textit{thermal\_1 變動成本低（基載）；thermal\_2 啟動便宜但每 MWh 較貴（尖峰備用）。兩部初始皆關機。}

\subsection{再生能源與儲能}
\begin{center}
\begin{tabularx}{\linewidth}{llX}
\toprule
\textbf{設備} & \textbf{型別} & \textbf{關鍵參數} \\
\midrule
pv\_1 & 太陽能 & 額定 60 MW，依日前預測比例出力 \\
pv\_2 & 太陽能 & 額定 80 MW，預測曲線每日三段日照尖峰（約 8--13、31--43、55--66 時）\\
battery\_1 & 儲能 & SOC 20--100 MWh，充/放 20 MW，初始 45 MWh \\
battery\_2 & 儲能 & SOC 10--60 MWh，充/放 15 MW，初始 25 MWh \\
\bottomrule
\end{tabularx}
\end{center}
\noindent\textit{每個電池對應一個充電需求（battery\_1\_chg / battery\_2\_chg），充電只能由機組或再生能源供應，不能由其他電池放電供應。}

\subsection{市場價格}
\fn{price\_72hr.json} 提供每小時售電價 $\lambda_t$。本資料尖峰落在第 70 時（113）、第 22 時（108）、第 56 時（105），谷底為第 12 時（0）。排程因此傾向把多餘電力留到高價時段售出。

% ---- 3 架構 ----
\section{系統架構與管線（Pipeline）}
整個系統由 \fn{src/main.py} 串接四個階段；Level 2 另以 \fn{run\_level2()} 串成「生成 $\to$ 靜態 L1 $\to$ 動態 L2 $\to$ 比較」。
\begin{code}
Phase 1  任務生成   generator/          -> output/task_set.json
Phase 2  日前排程   rt_scheduler/       -> output/schedule_result.json  (PuLP MILP)
Phase 3  接收測試   acceptance_tester   -> 在排程上標註接受/拒絕 (整合於 Phase 2 之後)
Phase 4  效能評估   evaluator/          -> output/evaluation_results.json
Level 2  動態排程   advanced_scheduler  -> output/*_dynamic.json
\end{code}
設計上採物件導向與單一職責：生成器把計算交給 \fn{FrameSizeCalculator}、驗證交給 \fn{TaskSetValidator}；排程器把建模交給 \fn{VppMilpFormulator}、結果解析交給 \fn{SchedulerResultExtractor}。資料模型以 Pydantic 定義，I/O 集中於 \fn{JsonIO}。

% ---- 4 Phase 1 ----
\section{Phase 1：週期任務集生成}
\fn{TaskSetGenerator} 以「生成 $\to$ 計算 frame size $\to$ 驗證」迴圈，隨機產生一組同時滿足所有作業限制的週期任務；不合格即重抽，直到通過 \fn{TaskSetValidator}。為保證較難的約束，生成策略刻意：
\begin{itemize}
  \item 先生成 $d=e$ 的任務（保證評分 1-6：$\ge 20\%$ 任務 $d_j = e_j$）。
  \item 再生成 2 個 $e=2$ 的非搶占任務（保證評分 1-7：$\ge 2$ 個 $e\neq 1$ 的非搶占任務）。
  \item 其餘任務隨機補足，並維持週期種類 $\ge 3$。
\end{itemize}
Frame size $f$ 取最小可行值，需同時滿足：
\begin{equation}
f \ge \max_j e_j,\qquad H \bmod f = 0,\qquad 2f - \gcd(f, p_j) \le d_j \quad \forall j .
\end{equation}

\subsection{本次產生的週期任務集（8 個任務，frame size $=4$）}
\begin{center}
\begin{tabularx}{\linewidth}{cccccccc}
\toprule
\textbf{ID} & \textbf{$r$ 釋放} & \textbf{$p$ 週期} & \textbf{$e$ 執行} & \textbf{$d$ 截止} & \textbf{$w$(MWh)} & \textbf{搶占} & \textbf{job 數} \\
\midrule
p1 & 20 & 20 & 4 & 4  & 11 & 否 & 3 \\
p2 & 1  & 10 & 2 & 8  & 17 & 否 & 7 \\
p3 & 17 & 20 & 4 & 4  & 8  & 否 & 3 \\
p4 & 8  & 8  & 1 & 5  & 13 & 否 & 8 \\
p5 & 8  & 24 & 2 & 19 & 17 & 否 & 2 \\
p6 & 11 & 13 & 2 & 11 & 17 & 否 & 4 \\
p7 & 10 & 16 & 2 & 13 & 16 & 否 & 4 \\
p8 & 9  & 9  & 2 & 7  & 13 & 否 & 7 \\
\bottomrule
\end{tabularx}
\end{center}
合計展開為 \textbf{38 個週期 job}（$>30$ \checkmark）；密度 $D_W = \sum_j e_j/p_j = 1.31 \ge 0.7$ \checkmark；週期 7 種（$\ge 3$ \checkmark）；$d=e$ 者為 p1、p3（$2/8 = 25\%\ge 20\%$ \checkmark）；非搶占且 $e\neq 1$ 者有 p2、p5、p6、p7、p8（$\ge 2$ \checkmark）；$f=4$ 通過上式（p8 最緊：$2f-\gcd(4,9)=7\le d_{\text{p8}}=7$）\checkmark。本次實例所有任務皆為非搶占。

\noindent 展開規則：第 $k$ 個 job 的 $\text{release}=r+kp$，$\text{deadline}=\text{release}+d-1$，僅保留 deadline $\le 72$ 者。此外 \fn{task\_set.json} 亦含 Demo 提供的 6 個零星任務（s1--s6）與 10 個非週期任務（a1--a10），供 Phase 3 使用。

% ---- 5 Phase 2 ----
\section{Phase 2：日前 MILP 排程}
\fn{VppMilpFormulator} 把 72 小時排程建成一個混合整數線性規劃（MILP），以 PuLP + CBC 一次解出整個 horizon。

\subsection{決策變數}
\begin{center}
\begin{tabularx}{\linewidth}{lX}
\toprule
\textbf{變數} & \textbf{意義} \\
\midrule
$P_{i,t}$ & 設備 $i$ 在 $t$ 的總出力（MWh） \\
$k_{j,i,t}$ & 設備 $i$ 在 $t$ 供給需求 $j$ 的能量（充電需求則為對電池的充電量） \\
$u_{i,t},\,\mathrm{start}_{i,t},\,\mathrm{stop}_{i,t}$ & 機組開關機／啟動／停機二元變數（線性化 $\min(1,P)$） \\
$\mathrm{charge\_b}_{i,t},\,\mathrm{discharge\_b}_{i,t}$ & 電池充／放電二元旗標（互斥） \\
$\mathrm{SOC}_{i,t}$ & 電池在 $t$ 末的儲能狀態 \\
$\mathrm{Sell}_t$ & $t$ 時售電量 \\
$x_{j,t}$ & 需求 $j$ 在 $t$ 是否執行之二元變數（線性化 $\min(1,\textstyle\sum k)$） \\
\bottomrule
\end{tabularx}
\end{center}
\noindent\textit{規格中的 $\min(1,\cdot)$ 為非線性，實作以二元變數 $x$（需求是否啟動）與 $u$（機組是否開機）線性化——這是把規格轉成可解 MILP 的關鍵。}

\subsection{核心限制式}
任務完成性與供需平衡是模型骨幹，以數學式列出幾條代表：
\begin{align}
\text{(C1) 啟動時收滿需求：}\quad & \textstyle\sum_{i\in I} k_{j,i,t} = w_j\, x_{j,t}, && \forall j,\ \forall t \\
\text{(C3) 截止前累積執行 $e$ 格：}\quad & \textstyle\sum_{t=r_j}^{r_j+d_j-1} x_{j,t} = e_j, && \forall j\in J_p\cup J_s \\
\text{(C7) 爬升率：}\quad & P_{i,t}-P_{i,t-1}\le RU_i\,\Delta t,\ \ P_{i,t-1}-P_{i,t}\le RD_i\,\Delta t, && \forall i\in I_g \\
\text{(C13) 再生能源上限：}\quad & P_{i,t}\le \mathrm{renewmax}_i\cdot\mathrm{forecast}_{i,t}\cdot\Delta t, && \forall i\in I_r \\
\text{(C16) SOC 平衡：}\quad & \mathrm{SOC}_{i,t}=\mathrm{SOC}_{i,t-1}+\!\!\sum_{\text{charge}}\!\! k - P_{i,t}, && \forall i\in I_b \\
\text{(C23) 供需平衡：}\quad & \textstyle\sum_{i\in I}P_{i,t}=\sum_{j}\sum_{i}k_{j,i,t}+\mathrm{Sell}_t, && \forall t
\end{align}

\subsection{23 條限制式總覽（依實作分組）}
\begin{center}
\begin{tabularx}{\linewidth}{llX}
\toprule
\textbf{群組} & \textbf{編號} & \textbf{內容} \\
\midrule
任務 & C1--C3, C5 & 啟動收滿 $w$；釋放前不執行；截止前累積 $e$ 格；非搶占須連續（$0\to1$ 上升沿 $\le 1$） \\
非週期 & C4 & $\mathrm{Miss}_j$ 定義：未逾時即完成 $e$ 格，否則記 $\mathrm{Miss}=1$ \\
機組 & C6--C12 & 出力上下限、爬升率、下限 $\le$ 單格爬升、最小開/關機、初始開/關機延續 \\
再生能源 & C13 & $P\le \mathrm{capacity}\cdot\mathrm{forecast}\cdot\Delta t$ \\
儲能 & C14--C19, C21 & 充/放上限、SOC 平衡、SOC 上下限、放電 $\le$ 可用儲能、不可同時充放、充電不可由非發電設備供應 \\
售電 & C22 & $\mathrm{Sell}_t\ge 0$ \\
平衡 & C20, C23 & 各設備供給 $\le$ 其出力；每小時總發電 $=$ 需求 $+$ 充電 $+$ 售電 \\
\bottomrule
\end{tabularx}
\end{center}
\noindent\textit{C16 與 C18 在 Level 2 由「真實儲能」放寬版取代（見第 8 節），預設參數下與原式等價。}

\subsection{目標函數與求解}
目標即第 1 節的 $F$。Level 1 的 MILP 只將週期 job 與充電 job 納入模型；零星／非週期任務由 Phase 3 以保留量處理，故靜態 MILP 求解時 $f_1=0$，非週期逾時罰則於 Phase 4 評估時才計入 \fn{objective\_value}。求得最佳解後，\fn{SchedulerResultExtractor} 解析各變數、四捨五入、捨去近零值，輸出 \fn{schedule\_result.json}。

% ---- 6 Phase 3 ----
\section{Phase 3：接收測試（Sporadic／Aperiodic）}
\fn{AcceptanceTester} 在 MILP 解出後執行（已整合進 \fn{RTScheduler.run()}）。它利用 MILP 留下的\textbf{保留量}——原本要賣到市場的剩餘電力 $\mathrm{Sell}_t$，可拆解為各設備的 spare：
\begin{equation}
\mathrm{spare}_{i,t}=P_{i,t}-\textstyle\sum_{j}k_{j,i,t},\qquad \sum_i \mathrm{spare}_{i,t}=\mathrm{Sell}_t\ (\text{由 C23}).
\end{equation}
接受 job 時，把其每格需求 $w$ 從 spare 轉入 $k_{\text{job},i,t}$，並同額減少 $\mathrm{Sell}_t$——總發電不變，C23、C20 仍成立。決策規則：
\begin{itemize}
  \item \textbf{零星（硬截止）}：唯有 $[r,\,r+d-1]$ 內能排入 $e$ 格、每格保留量 $\ge w$ 時才接受，否則拒絕。
  \item \textbf{非週期（軟截止）}：盡量排到 $H$ 前完成；若完成時刻晚於軟截止則另記為 miss。
  \item 搶占 $=1$ 可用不連續時槽；搶占 $=0$ 須單一連續視窗。
\end{itemize}
每個決策（接受/拒絕、排入時槽、完成時刻、理由）寫入 \fn{acceptance\_test\_log.json}，支援評分 4-1（方法）與 4-2（合理性）。

\subsection{本次靜態接收測試結果}
零星 s1--s6 全部接受（value rate $=1.0$）：例如 s1 在 $[2,9]$ 排入第 2、4 時；s3 在 $[28,33]$ 排入第 28、29 時。非週期 a1--a10 全部排入，但 a2（完成於 21 $>$ 軟截止 17）與 a5（完成於 46 $>$ 軟截止 42）逾時，故靜態軟逾時率 $=2/10=0.2$。

% ---- 7 Phase 4 ----
\section{Phase 4：效能評估指標}
\fn{Evaluator} 讀入排程結果與輸入檔，計算所有規格指標，輸出 \fn{evaluation\_results.json}：
\begin{center}
\begin{tabularx}{\linewidth}{lX}
\toprule
\textbf{指標} & \textbf{定義} \\
\midrule
hard\_deadline\_miss\_rate & （逾時的週期 $+$ 零星 job）$/$ 總硬截止 job 數 \\
soft\_deadline\_miss\_rate & 逾時的非週期 job $/$ 總非週期 job 數 \\
average / max\_tardiness & $T_j=\max(0,\,C_j-d_j)$ 的平均與最大 \\
average / max\_response\_time & $R_j=C_j-r_j$ 的平均與最大 \\
completion\_time\_jitter & 同一週期任務各 instance 完成時刻的峰對峰值，再取平均 \\
sporadic\_value\_rate & 硬截止前完成的零星執行格數 $/$ 零星總執行格數 \\
generator\_cost / market\_revenue & $f_2$ 機組成本 $/$ $\sum_t \lambda_t\,\mathrm{Sell}_t$ 售電收益 \\
objective\_value & $F=\alpha\cdot(\text{軟逾時數})+\text{機組成本}-\text{售電收益}$ \\
\bottomrule
\end{tabularx}
\end{center}

% ---- 8 Level 2-A ----
\section{Level 2-A：放寬假設建模（R1--R10）}
我們放寬三個 Level 1 假設——假設 11（再生能源預測必準）、假設 12（理想儲能）、假設 5（無工作優先序）——並嚴格遵守規格：\textbf{可修改符號、可加參數，但不得新增決策變數}。下列每條放寬限制都只用 Level 1 既有變數 $P,\,k,\,\mathrm{SOC},\,\mathrm{Sell},\,x$ 表達。所有參數集中於 \fn{runtime\_config.json}；全取預設（$\eta=1,\sigma=0,\beta=0,\dots$）時模型與 Level 1 完全相同。

\subsection{放寬參數表（本次採用值）}
\begin{center}
\begin{tabularx}{\linewidth}{llccX}
\toprule
\textbf{參數} & \textbf{符號} & \textbf{值} & \textbf{假設} & \textbf{意義} \\
\midrule
charge\_efficiency & $\eta_c$ & 0.95 & 12 & 充電能量僅 $\eta_c$ 進入電芯 \\
discharge\_efficiency & $\eta_d$ & 0.95 & 12 & 放出 $P$ 需從電芯抽 $P/\eta_d$ \\
self\_discharge\_rate & $\sigma$ & 0.002 & 12 & 每格自放電比例 \\
cycle\_limit & $N$ & 3.0 & 12 & 全期最大等效循環數 \\
soc\_power\_floor & $\varphi$ & 0.5 & 12 & SOC 極端時最低功率比例 \\
aging\_cost & $c_{\text{age}}$ & 2.0 & 12 & 每放電 1 MWh 老化成本(\$) \\
renewable\_uncertainty\_margin & $\beta$ & 0.15 & 11 & 預測折減保留邊際 \\
precedence & --- & 自動 & 5 & 有序工作對 $(a,b)$ \\
\bottomrule
\end{tabularx}
\end{center}

\subsection{放寬限制式 R1--R10}
\textbf{假設 11 —— 再生能源不確定性}
\begin{align}
\text{(R1 / C13$'$)}\quad & P_{i,t}\le \mathrm{renewmax}_i\cdot\mathrm{forecast}_{i,t}\cdot(1-\beta)\cdot\Delta t, && \forall i\in I_r .
\end{align}
日前計畫不靠滿預測，保留邊際以防高估；$\beta=0$ 還原 C13。

\noindent\textbf{假設 12 —— 真實儲能（就地改寫 C16、C18，預設參數即還原）}
\begin{align}
\text{(R2--R4 / C16$'$)}\quad & \mathrm{SOC}_{i,t}=(1-\sigma)\,\mathrm{SOC}_{i,t-1}+\eta_c\,\mathrm{charge\_in}_{i,t}-\tfrac{1}{\eta_d}P_{i,t} \\
\text{(R5 / C18$'$)}\quad & \tfrac{1}{\eta_d}P_{i,t}\le (1-\sigma)\,\mathrm{SOC}_{i,t-1}-\mathrm{soc\_min}_i \\
\text{(R6 循環/吞吐)}\quad & \textstyle\sum_{t}P_{i,t}\le N\,(\mathrm{soc\_max}_i-\mathrm{soc\_min}_i) \\
\text{(R7 SOC 相依放電)}\quad & P_{i,t}\le \mathrm{dis}_i\,\Delta t\Big[\varphi+(1-\varphi)\tfrac{\mathrm{SOC}_{i,t-1}-\mathrm{soc\_min}_i}{\mathrm{soc\_max}_i-\mathrm{soc\_min}_i}\Big] \\
\text{(R8 SOC 相依充電)}\quad & \mathrm{charge\_in}_{i,t}\le \mathrm{chg}_i\,\Delta t\Big[\varphi+(1-\varphi)\tfrac{\mathrm{soc\_max}_i-\mathrm{SOC}_{i,t-1}}{\mathrm{soc\_max}_i-\mathrm{soc\_min}_i}\Big]
\end{align}
\begin{align}
\text{(R9 老化成本，進目標)}\quad & F \mathrel{+}= \sum_{i\in I_b}\sum_{t}c_{\text{age}}\,P_{i,t} .
\end{align}
R7/R8 使功率在接近滿電時達額定、接近不利端時降到 $\varphi\cdot$額定，線性內插。

\noindent\textbf{假設 5 —— 工作優先序}
\begin{align}
\text{(R10)}\quad & e_a\, x_{b,t}\le \sum_{s<t}x_{a,s}, && \forall t\in b\text{ 的視窗}.
\end{align}
即 $a$ 全部做完前 $b$ 不得啟動；重用 $x$，無新變數。未設定時程式自動挑一組可行的非平凡對（本次為 \fn{p2\_0} $\to$ \fn{p4\_0}）。

\noindent\textbf{保留策略（支援評分 6/7 與動態法）}
\begin{align}
\text{(R-reserve)}\quad & \mathrm{Sell}_t\ge R_t .
\end{align}
以對 $\mathrm{Sell}$ 設下限強制計畫保留可轉移剩餘，供接收測試使用；重用 $\mathrm{Sell}$，無新變數。

\medskip
\noindent 合計 \textbf{10 條放寬限制（R1--R10）} 外加保留下限，達到評分 3 的上限 10 分，且皆未新增決策變數。

% ---- 9 Level 2-B ----
\section{Level 2-B：進階動態排程（滾動視窗再最佳化）}
\textbf{方法：Receding-horizon（滾動視窗）再最佳化。} Level 1 在 $t=0$ 一次解完整個 $[1,72]$；\fn{AdvancedScheduler} 則沿時間前進，在觸發點對「剩餘 horizon」重新最佳化，因應靜態計畫看不到的動態資訊。

\subsection{核心機制}
\begin{itemize}
  \item \textbf{狀態承接（pin\_prefix）}：每次觸發時，已執行的 $[1,t_0)$ 被釘住為已承諾值（連續變數 $P,k,\mathrm{SOC},\mathrm{Sell}$ 加上由其推導的前綴二元變數 $u,\mathrm{charge\_b},\mathrm{discharge\_b},x$），只重解 $[t_0,72]$。釘住前綴讓 solver 的 presolve 將其消去，使每次再解都便宜，且機組開關機時長、爬升位置、SOC 都正確跨界承接。
  \item \textbf{觸發時機}：(a) 週期節奏 \fn{reopt\_interval}$=24$ 格；(b) 每個零星任務釋放（硬任務需即時審核）；(c) 最多 3 個「實際 PV 與預測偏差最大」的時點。非週期任務於下一個節奏邊界揭露。
  \item \textbf{再生能源實現}：以種子化隨機模型（noise std $=0.2$, seed $=42$）在預測周圍生成「實際」PV。已承諾區塊用實際值、未見尾段用折減預測（R1）——對已承諾精確、對未來保守。
  \item \textbf{線上接收 $=$ 保留可行性}：任務於釋放時揭露，能否接受取決於再最佳化「能否在其執行視窗內保留足夠可轉移剩餘（$\mathrm{Sell}\ge R_t$）」。硬任務不可行即拒絕，軟任務盡力保留。被接受者隨區塊凍結被路由進排程，故接受的硬任務必在截止前完成。
\end{itemize}
\noindent\textit{每次再解使用 MIP gap（\fn{gap\_rel}$=0.05$）與時間上限（\fn{time\_limit}$=20$s），確保線上再最佳化可解。}

\subsection{目標衝突（為何無法同時最佳化三目標）}
\begin{itemize}
  \item 提高接收（降 miss）$\Rightarrow$ 須保留更多剩餘 $\Rightarrow$ 更多機組承諾 $\Rightarrow$ 機組成本 $f_2$ 上升，或售電收益下降（剩餘被留著而非在尖峰賣出）。
  \item 再生能源不確定（R1 折減 $+$ 實際短缺）$\Rightarrow$ 能量由 PV 移到火力 $\Rightarrow$ $f_2$ 上升。
  \item 儲能真實化（效率、自放電、循環、老化）$\Rightarrow$ 電池成為更貴、更有損的緩衝 $\Rightarrow$ 套利收益變少。
\end{itemize}
因此動態法是「以較高機組成本換取更少逾時與更佳即時回應」，下一節以實際數據驗證。

% ---- 10 結果 ----
\section{實際執行結果與靜態 vs 動態比較}
\subsection{動態排程的 9 輪再最佳化（共 9 次 MILP 解，約 12 秒，單執行緒、可重現）}
\begin{center}
\begin{tabularx}{\linewidth}{cclX}
\toprule
\textbf{觸發 $t_0$} & \textbf{承諾區間} & \textbf{本輪揭露} & \textbf{結果} \\
\midrule
1  & $[1,1]$   & （無）        & 建立初始計畫 \\
2  & $[2,13]$  & s1, a1, a2   & 全部接受並路由 \\
14 & $[14,24]$ & s2, a3       & 全部接受 \\
25 & $[25,27]$ & a4           & 接受 \\
28 & $[28,39]$ & s3, a5       & 全部接受 \\
40 & $[40,48]$ & s4, a6       & 全部接受 \\
49 & $[49,52]$ & a7           & 接受 \\
53 & $[53,65]$ & s5, a8, a9   & 全部接受 \\
66 & $[66,72]$ & s6, a10      & 全部接受 \\
\bottomrule
\end{tabularx}
\end{center}
9 輪皆無拒絕、無再排程失敗；38 個週期 job、6 個零星 job 全部於硬截止前完成，10 個非週期 job 全部於軟截止前完成。逐格驗證：C23 平衡 0 違反、SOC 維持於 $[\mathrm{soc\_min},\mathrm{soc\_max}]$ 0 違反、放寬儲能/機組/再生能源限制每輪皆成立。

\subsection{靜態（L1）vs 動態（L2）效能比較}
\begin{center}
\begin{tabularx}{\linewidth}{Xrrr}
\toprule
\textbf{指標} & \textbf{靜態 L1} & \textbf{動態 L2} & \textbf{差異 $\Delta$} \\
\midrule
objective\_value & $-69\,410.4$ & $\mathbf{-83\,238.8}$ & $-13\,828.4$（更佳） \\
generator\_cost & $296\,403.6$ & $311\,428.9$ & $+15\,025.3$ \\
market\_revenue & $385\,814.0$ & $394\,667.8$ & $+8\,853.8$ \\
hard\_deadline\_miss\_rate & $0.0$ & $0.0$ & --- \\
soft\_deadline\_miss\_rate & $\mathbf{0.2}$ & $\mathbf{0.0}$ & $-0.2$ \\
average\_tardiness & $0.15$ & $0.0$ & $-0.15$ \\
max\_tardiness & $4$ & $0$ & $-4$ \\
average\_response\_time & $3.63$ & $2.65$ & $-0.98$ \\
sporadic\_value\_rate & $1.0$ & $1.0$ & --- \\
\bottomrule
\end{tabularx}
\end{center}
\noindent\textbf{解讀：}
\begin{itemize}
  \item 靜態日前計畫雖一次看到所有任務，但只承諾一份排程，漏掉 2 個非週期任務（a2、a5，軟逾時率 0.2）。動態法隨任務到達再最佳化並為其保留容量，10 個全部準時（軟逾時率 0.0），平均回應時間由 3.63 降到 2.65。
  \item 消除 2 個軟逾時 $=$ 省下 $2\alpha=20\,000$ 罰則，主導目標改善（$-69\,410 \to -83\,239$）。
  \item 代價清晰可見：保留容量並補償實際再生能源短缺使機組成本上升（$+15\,025$）；較大的承諾發電同時讓售電收益增加（$+8\,854$）。扣掉逾時罰則後，動態法在總目標上明顯更佳，印證上節的權衡分析。
\end{itemize}

% ---- 11 程式結構 ----
\section{程式碼結構與執行方式}
\begin{code}
src/
  main.py              管線進入點 (run_level2 = L1 + L2 + 比較)
  config.py            路徑與常數
  validator.py         自我檢核 (純標準庫)
  advanced_scheduler.py  Level 2 滾動視窗動態排程
  generator/           Phase 1 任務生成
  rt_scheduler/        Phase 2+3: formulator / acceptance_tester /
                        relaxation / expander / extractor
  evaluator/           Phase 4 評估
  model/               Pydantic 資料模型
input/   processor_settings.json, price_72hr.json
output/  task_set.json, schedule_result*.json, evaluation_results*.json,
         acceptance_test_log*.json, dynamic_run_log.json
runtime_config.json    Level 2 放寬 + 動態節奏參數
\end{code}
執行指令：
\begin{code}
# 環境
pip install pulp pydantic          # 或 poetry install

# Level 1 完整管線 (生成 -> 排程 -> 評估)
python -m src.main

# Level 2 比較管線 (生成 -> 靜態 -> 動態 -> 印出比較)
python -c "from src.main import run_level2; run_level2()"

# 只跑動態排程 / 自我檢核
python -m src.advanced_scheduler
python3 -m src.validator              # Level 1
python3 -m src.validator --level 2    # Level 2
\end{code}

% ---- 12 自我檢核 ----
\section{對應評分表的自我檢核}
\begin{center}
\begin{tabularx}{\linewidth}{Xcl}
\toprule
\textbf{評分項} & \textbf{配分} & \textbf{本作業狀態} \\
\midrule
1 週期任務集（1-1$\sim$1-8） & 17 & 全數通過（見 §4.1） \\
2 模型限制（C1$\sim$C23，含放寬版） & 27 & 全數實作並通過自我檢核 \\
3 放寬假設（R1--R10） & 10 & 10 條皆由動態排程實際滿足 \\
4 接收測試（含 4-3 value rate） & 11 & 零星 value rate $=1.0$；附逐 job 理由 \\
5 排程結果 & 8 & JSON 格式相容；週期 job 全準時 \\
6 評估指標 & 7 & 全部指標皆輸出 \\
7 保留策略分析 & 10 & 報告分析 $+$ 保留下限 R-reserve \\
8 動態法（8-1/8-2/8-3） & 10 & 設計 $+$ 正確性 $+$ 靜/動比較皆具備 \\
\bottomrule
\end{tabularx}
\end{center}
\noindent\textit{自我檢核工具 \fn{python3 -m src.validator -{}-level 2} 會用放寬後的儲能/再生能源限制重驗每一項，並逐條確認動態排程確實滿足 R1--R10（例如 SOC 是否遵循效率/自放電平衡、再生能源是否守住實際上限、優先序是否成立），亦即驗證「實作」而非僅「描述」。}

% =========================================================== 第二部分
\newpage
\part*{第二部分　報告講稿（明日對老師報告）}
\addcontentsline{toc}{section}{第二部分　報告講稿}
\textit{建議總時長 10--12 分鐘。以下每段對應一張投影片，粗體為口語講稿，項目為提醒重點。}

\subsection*{投影片 1：開場與題目定位（約 1 分鐘）}
\textbf{講稿：}老師好，我們這組做的是「虛擬電廠的即時排程系統」。核心想法是：把電廠裡的傳統機組、太陽能、電池當成即時系統的『處理器』，把各種有時間限制的用電需求當成『任務』，於是供電調度就變成一個即時排程問題。我們要在 72 小時、每小時一格的時間軸上安排這些處理器供電，同時滿足三個彼此衝突的目標：少漏非週期任務、降低機組成本、提高售電收益。
\begin{itemize}
  \item 先讓老師建立「處理器 $=$ 發電設備、任務 $=$ 用電需求」的對應。
  \item 強調三個目標互相衝突，這是後面 Level 2 分析的伏筆。
\end{itemize}

\subsection*{投影片 2：系統架構與四階段管線（約 1 分鐘）}
\textbf{講稿：}整個系統分四個階段：Phase 1 生成週期任務集；Phase 2 用混合整數線性規劃解出日前排程；Phase 3 在排程留下的保留量上做接收測試，處理零星與非週期任務；Phase 4 計算所有效能指標。Level 2 我們在這個基礎上放寬三個假設，並加上一個會隨資訊滾動再最佳化的動態排程器。
\begin{itemize}
  \item 指著架構圖由左到右講一次資料流。
  \item 點出 Phase 3 已整合進排程器之後自動執行。
\end{itemize}

\subsection*{投影片 3：Phase 1 週期任務生成（約 1 分鐘）}
\textbf{講稿：}週期任務是用「生成—驗證」迴圈產生的，會反覆重抽直到同時滿足作業所有限制；我們加了固定種子讓它可重現。我們刻意先生成 $d=e$ 的任務和非搶占任務來保證最難滿足的條件。這次產生 8 個週期任務，frame size 是 4，展開後是 38 個 job，超過 30 個；密度 1.31 遠高於 0.7 的下限；frame size 也通過 $2f-\gcd(f,p)\le d$ 這條。
\begin{itemize}
  \item 報出關鍵數字：8 個任務、frame$=4$、38 個 job、密度 1.31。
  \item 若問為何是 frame size，說明它保證每個 job 有完整一格可被排進去。
\end{itemize}

\subsection*{投影片 4：Phase 2 MILP 模型——變數與目標（約 1.5 分鐘）}
\textbf{講稿：}日前排程是一個 MILP，我們用 PuLP 加 CBC 一次解出整個 72 小時。決策變數有設備出力 $P$、能量分配 $k$、機組開關機、電池充放電與 SOC、售電量 Sell，還有需求是否執行的二元變數 $x$。目標函數是 $\alpha$ 乘上非週期逾時數，加機組成本，減售電收益；$\alpha$ 設成 10000，代表每漏一個非週期任務罰一萬元。規格裡的 $\min(1,\cdot)$ 是非線性，我們用二元變數 $x$ 和 $u$ 把它線性化，這樣才能交給 MILP 解。
\begin{itemize}
  \item 強調「一次解完整個 horizon」是靜態法的特徵，與 Level 2 對照。
  \item 解釋 $\alpha=10000$ 把「個數」換算成「金額」才能和成本相加。
\end{itemize}

\subsection*{投影片 5：Phase 2 模型——23 條限制式（約 1.5 分鐘）}
\textbf{講稿：}限制式共 23 條，分成幾群：任務類保證每個 job 在截止前累積執行恰好 $e$ 格、非搶占要連續；機組類包含出力上下限、爬升率、最小開關機時間；再生能源受預測上限約束；儲能類包含充放電上限、SOC 平衡、不可同時充放；最後是每小時的供需平衡 C23——每一格的總發電必須等於用電需求加充電加售電。這條平衡式是整個模型的骨幹。
\begin{itemize}
  \item 不必逐條念，挑 C23 平衡與 C3 完成性、C5 連續性講透即可。
  \item 提一句 C16/C18 在 Level 2 會被放寬版取代。
\end{itemize}

\subsection*{投影片 6：Phase 3 接收測試與保留策略（約 1.5 分鐘）}
\textbf{講稿：}排程解完後，市場上沒賣掉的剩餘電力就是我們的「保留量」。接收測試就靠這個保留量來接零星和非週期任務。零星任務是硬截止：只有當它的時間視窗內每一格的保留量都夠它的能量需求、而且能排滿執行格數時，才接受，否則拒絕。非週期任務是軟截止：盡量排，逾時就記成 miss。接受時我們把需求從設備的剩餘出力轉進去、同時把售電量等額扣掉，所以總發電不變、平衡式仍然成立。這次 6 個零星任務全部接受，value rate 是 1.0。
\begin{itemize}
  \item 講清楚「轉移剩餘、扣售電、平衡不變」這個關鍵手法。
  \item 報出：零星全接受、value rate 1.0；靜態下 a2、a5 兩個非週期逾時。
\end{itemize}

\subsection*{投影片 7：Level 2-A 放寬三個假設（約 1.5 分鐘）}
\textbf{講稿：}Level 2 我們放寬三個假設：再生能源預測不一定準、儲能不是理想的、以及工作之間有優先序。最重要的限制是：規格規定可以加參數、改符號，但不能新增決策變數，所以我們所有放寬限制都只用原本的變數來寫。再生能源我們把預測折減 15\% 留邊際；儲能我們加了充放電效率 0.95、自放電、循環上限、SOC 相依功率和老化成本；優先序則是強制某個 job 做完另一個才能開始。總共 10 條放寬限制，剛好對到評分的 10 分上限。把參數設回預設，模型就跟 Level 1 完全一樣。
\begin{itemize}
  \item 反覆強調「沒有新增決策變數」——這是評分硬規定。
  \item 舉一兩個具體：SOC 平衡加了效率與自放電；老化成本進目標。
\end{itemize}

\subsection*{投影片 8：Level 2-B 動態排程方法（約 1.5 分鐘）}
\textbf{講稿：}動態排程用的是滾動視窗再最佳化。Level 1 是一次解完全程，動態法則是沿時間往前走，在三種觸發點重新最佳化剩下的時程：固定每 24 小時一次、每個零星任務一到就審核、還有實際太陽能跟預測差最大的幾個點。已經執行過的部分會被「釘住」，只重解後面，這樣每次再解都很快，而且機組狀態跟 SOC 都正確接續。已承諾的區塊用「實際」的太陽能出力、看不到的未來用折減預測，所以對已承諾精確、對未來保守。任務能不能接受，看再最佳化能不能在它的視窗內保留足夠剩餘，這就是把接收測試做成 MILP 的可行性問題。
\begin{itemize}
  \item 強調「釘住前綴 $+$ 只重解尾段」是讓線上再最佳化可解的關鍵。
  \item 用一句話帶出：硬任務接受了就保證準時完成。
\end{itemize}

\subsection*{投影片 9：結果——9 輪再最佳化與正確性（約 1 分鐘）}
\textbf{講稿：}這次動態排程跑了 9 輪、9 次 MILP，大約 12 秒（單執行緒、可重現）。任務隨時間陸續到達並被接受：例如第 2 小時揭露 s1、a1、a2，第 28 小時揭露 s3、a5，一路到第 66 小時的 s6、a10。最後所有 38 個週期 job、6 個零星 job 都在硬截止前完成，10 個非週期 job 也都在軟截止前完成；硬逾時率 0、軟逾時率 0。每一格的供需平衡和 SOC 範圍我們也都逐格驗證沒有違反。
\begin{itemize}
  \item 報出：9 輪、9 次解、約 12 秒、全部準時。
  \item 強調「逐格驗證 0 違反」回應評分 8-2 的正確性要求。
\end{itemize}

\subsection*{投影片 10：結果——靜態 vs 動態比較（約 1.5 分鐘）}
\textbf{講稿：}最關鍵的比較在這裡。靜態法雖然一次看到所有任務，但只承諾一份計畫，結果漏掉 2 個非週期任務，軟逾時率 0.2。動態法隨任務到達再最佳化、並主動保留容量，10 個非週期任務全部準時，軟逾時率變成 0，平均回應時間也從 3.63 降到 2.65。少漏 2 個就省下 2 萬元罰則，這主導了總目標從 $-69410$ 改善到 $-83239$。代價是：保留容量加上補償太陽能實際短缺，機組成本多了約 1.5 萬；不過較大的發電也讓售電收益增加約 8.9 千。這正好驗證我們前面說的——提高接收率要用機組成本來換，三個目標沒辦法同時最佳化。
\begin{itemize}
  \item 這是全場重點，務必把「省 2 萬罰則 $>$ 多花 9 千成本」的帳算給老師聽。
  \item 結論句：動態法在總目標上明顯更好，且權衡關係與理論分析一致。
\end{itemize}

\subsection*{投影片 11：結語與分工／AI 協作（約 1 分鐘）}
\textbf{講稿：}總結：我們完成了 Level 1 的生成、日前排程、接收測試與評估，也完成 Level 2 的三項假設放寬與動態排程，並用實際數據證明動態法在不漏硬截止的前提下，把軟逾時降到 0、總目標更好。程式採物件導向、單一職責設計，所有結果都可用固定種子重現。謝謝老師，以下開放提問。
\begin{itemize}
  \item 若需報告 AI 協作：說明用 AI 輔助建模與程式骨架、自己負責驗證與調參，並交叉檢查每條限制式。
  \item 預留時間回答提問。
\end{itemize}

\subsection*{預期問答（備用）}
\begin{itemize}
  \item \textbf{Q：為何用 MILP 而不是啟發式？}　A：MILP 能在 23 條硬限制下保證可行且最佳，且線上再最佳化用 gap 與時間上限即可在秒級內解出。
  \item \textbf{Q：動態法的計算成本？}　A：本例 9 次解約 12 秒（單執行緒）；釘住前綴讓 presolve 消去已承諾段，再解規模隨時間縮小。
  \item \textbf{Q：保證沒有新增決策變數？}　A：所有放寬限制（R1--R10）與保留下限都只用 $P,k,\mathrm{SOC},\mathrm{Sell},x$，已用自我檢核逐條驗證。
  \item \textbf{Q：jitter 動態略增（$44.1\to 45.8$）為何仍較好？}　A：jitter 是同任務各 instance 完成時刻的離散度，動態法為保留容量微調了部分週期 job 的完成時刻，但硬截止全數守住、總目標更佳，屬可接受的取捨。
\end{itemize}

\vspace{1cm}
\begin{center}\textit{（報告結束）}\end{center}

\end{document}
